\documentclass[12pt]{article}

% LuaLaTeX: fontspec + polyglossia для корректной работы с кириллицей
\usepackage{fontspec}
\defaultfontfeatures{Ligatures=TeX,Scale=MatchLowercase}
\setmainfont{Times New Roman}
\newfontfamily\cyrillicfont{Times New Roman}
\usepackage{polyglossia}
\setmainlanguage{russian}

\usepackage[style=numeric, backend=biber]{biblatex} % Обязательный пакет для корректной сборки
\usepackage{graphicx}        						% Графика
\usepackage{microtype}       						% Микротипографика
\usepackage{amssymb}        						% Математические символы
\usepackage{amsmath}        						% Математика
\usepackage{tikz}           						% Рисование
\usetikzlibrary{arrows.meta}
\usepackage{pgfplots}       						% Графики
\pgfplotsset{compat=1.18}
\usepackage{circuitikz}     						% Электрические схемы
\usepackage{fix-cm}         						% Чтобы работал большой размер шрифта
\usepackage{anyfontsize}    						% Для любого размера шрифта
\usepackage{fancyhdr}        						% Колонтитулы

\newfontfamily\titletimesnewroman[Scale=1.4]{Times New Roman}
\newfontfamily\regulartimesnewroman[Scale=1.2]{Times New Roman}

% Настройка колонтитула для титульной страницы
\fancypagestyle{titlepage}{%
  \fancyhf{}%
  \fancyfoot[C]{\linespread{1.4}\selectfont\regulartimesnewroman Москва\par 2025\par}% центр футера: Москва и новая строка 2025
  \renewcommand{\headrulewidth}{0pt}% убрать линию вверху
  \renewcommand{\footrulewidth}{0pt}% убрать линию внизу
}

\usepackage[a4paper,left=20mm,right=20mm,top=20mm,bottom=20mm]{geometry}

\begin{document}
\hbadness=10000 % отключить предупреждения underfull

\begin{center}
	
Федеральное государственное бюджетное образовательное учреждение высшего\\
профессионального образования\\
«Национальный исследовательский институт МИЭТ»
\vskip2cm
кафедра электротехники
\vskip2cm
% локально убрать отступ первой строки (parindent=0) внутри minipage, шрифт сохраняется
{
  \setlength{\parindent}{0pt}
  \linespread{1.5}\selectfont % увеличить межстрочный интервал
  % первая строка — центр
  \centering\titletimesnewroman\bfseries
  Пояснительная записка к курсовому проекту\par
  % следующие строки — растянуть на всю ширину страницы
  \titletimesnewroman\bfseries
  \makebox[\textwidth][s]{на тему: «Линейные электрические цепи постоянного и}\par
  \makebox[\textwidth][s]{переменного тока. (Методы расчета). Переходные процессы.»}\par
}
\end{center}

\vskip8cm

{
  \linespread{1.1}\selectfont

  \noindent Выполнил : Андреев И.А. \par
  \noindent Группа №~ ЭН-25 \par
  \vskip0.5cm
  \noindent Преподаватель : Самохин Д.В.\par
}

% Временно уменьшить отступ футера, чтобы поднять текст футера вверх на титульной странице
\addtolength{\footskip}{-6mm}% уменьшить на 6mm
\thispagestyle{titlepage}

\clearpage
% Восстановить отступ футера для последующих страниц
\addtolength{\footskip}{6mm}% вернуть прежнее значение

\section*{\centering Вариант 1}

\section{Найти $U_C(t)$}

\begin{tabular}{l c}
  % Документация:
  % \begin{minipage}[t]{0.25\linewidth}
  % [t] — опция вертикального выравнивания: выравнивать верхнюю границу minipage по базовой линии окружения (top).
  % Альтернативы: [c] — по центру, [b] — по нижней границе.
  % {0.25\linewidth} — ширина minipage (25% от ширины строки).
  % Содержимое — строки в математическом режиме; '\\' делает перенос строки внутри minipage.
  % Символ '&' указывает на разделение колонок (minipage, вероятно, используется внутри табличной/выравнивающей среды).
  % \begin{tikzpicture}[scale=1,baseline=(current bounding box.center), european resistors]
  % — начало рисунка TikZ: scale задаёт масштаб, baseline выравнивает рисунок по указанной точке, 'european resistors' — стиль/библиотека для элементного оформления схем.
  \scalebox{1.2}{
    \begin{minipage}[c]{0.2\linewidth}
      \vspace{0.5cm}
      $E=30\,\text{В}$ \\ 
      $R_1=5\,\Omega$\\
      $R_2=2.5\,\Omega$ \\ 
      $R_3=2.5\,\Omega$ \\ 
      $R_4=2.5\,\Omega$ \\ 
      $C=4\cdot 10^{-5}$\,\text{Ф} \\
    \end{minipage}
  } & \begin{tikzpicture}[scale=1,baseline=(current bounding box.center), european resistors, american inductors]
    \draw (-0.5,0) to[voltage source, l=$E$] (-0.5,2) -- (-0.5,3)
    to[short, -] (0,3) to[R, l=$R_1$] (2,3) to[short, -*] (2.5,3)
    -- (5.5,3) to[R, l=$R_3$] (5.5,1) 
    to[R, l=$R_4$] (5.5,-1) to[short, -*] (2.5,-1) -- (-0.5, -1) -- (-0.5,0);
    \draw (2.5,3) to[R, l=$R_2$] (2.5,1) to[C, l=$C$] (2.5,-1);
    \draw (5.5,1) -- (7,1) -- (7,0.4);
    \draw (5.5,-1) -- (7,-1) to[short, -*] (7,-0.4) -- (7.1,0.4);
    \draw[fill=white] (7,0.4) circle (1.5pt);
    \draw[-{stealth}] (6.7,0.1) arc (120:50:0.7);
  \end{tikzpicture} \\
\end{tabular}

\subsection{Решение}

\textbf{1. Начальное условие до коммутации (при $t < 0$, ключ замкнут)}

В начальный момент ключ замкнут, ток идет через $R_1$, $R_3$ последовательно. Ветвь с $R_2$ и $C$ параллельна резистору $R_3$.
В установившемся режиме через конденсатор ток не идет $i_C = 0$, следовательно, ток через $R_2$ также равен нулю. Весь ток идет через $E \to R_1 \to R_3\to E$.
Напряжение на конденсаторе равно напряжению на параллельной ветви с резистором $R_3$.
\[
u_C(0_-) = \frac{E}{R_1 + R_3} \cdot R_3 = 4 \cdot 2.5 = 10\,\text{В}
\]

Напряжение на конденсаторе не может измениться скачком, следовательно, согласно II закону коммутации
\[
u_C(0_+) = u_C(0_-) = 10\,\text{В}
\]

\textbf{2. Установившийся режим после коммутации (при $t \to \infty$, ключ разомкнут)}

После размыкания ключа, ток идет через $E\to R_1 \to R_3\to R_4 \to E$. Ветвь с $R_2$ и $C$ теперь резисторам $R_3$ и $R_4$.
В установившемся режиме через конденсатор ток не идет $i_C = 0$, следовательно, ток через $R_2$ также равен нулю.
Напряжение на конденсаторе равно напряжению на $R_3$ и $R_4$.
\[
u_{C\text{ пр}_+} = \frac{E}{R_1 + R_3 + R_4} \cdot (R_3 + R_4) = 3 \cdot 5 = 15\,\text{В}
\]

\textbf{3. Постоянная времени $\tau$ и показатель $p$}

Для определения $\tau$ необходимо найти эквивалентное сопротивление относительно конденсатора при закороченном источнике ЭДС и замкнутом ключе.

\[
R_{\text{экв}} = R_2 + \frac{R_1 \cdot R_3}{R_1 + R_3} = 2.5 + \frac{5 \cdot 2.5}{7.5} = 2.5 + \frac{12.5}{7.5} = 2.5 + \frac{5}{3} = \frac{25}{6} \approx 4.17\,\Omega
\]

Постоянная времени
\[
\tau = R_{\text{экв}} \cdot C = \frac{25}{6} \cdot 4 \cdot 10^{-5} = \frac{100 \cdot 10^{-5}}{6} = \frac{10^{-3}}{6}\,\text{с} \approx 0{,}167\,\text{мс}
\]

В свою очередь, показатель 
\[
p = -\frac{1}{\tau} = -\frac{6}{10^{-3}} = -6000\,\text{с}^{-1}
\]

\textbf{4. Свободная компонента}

Переходный процесс является апериодическим, так как в цепи присутствует только один энергонакопительный элемент. Чьи колебания описываются свободной компонентой напряжения.
\[
u_{C\text{ св}}(t) = A \cdot e^{pt}
\]

\textbf{5. Поиск постоянной интегрированния $A$}

Переходный процесс описывается уравнением
\[
\begin{aligned}
u_C(0_+) &= u_{C\text{ пр}}(0_+) + u_{C\text{ св}}(0_+) \\
10 &= 15 + A \cdot e^{0} \\
A &= 10 - 15 = -5\,\text{В}
\end{aligned}
\]
\textbf{6. Итоговое выражение для $u_C(t)$}
\[
u_C(t) = 15 - 5 \cdot e^{-6000t}\,\text{В}
\]

\begin{center}
\begin{tikzpicture}[xscale=12, yscale=0.75]
  % Сетка
  \draw[gray!30, dashed, ystep=1, xstep=0.1] (0,8) grid (0.99,15.9);
  
  % Оси
  \draw[-{stealth}, thick] (0,8) -- (1,8) node[right] {$t$, мс};
  \draw[-{stealth}, thick] (0,8) -- (0,16) node[above] {$u_C(t)$, В};
  
  % Деления на оси X (каждые 0.2 мс)
  \foreach \x in {0.2, 0.4, 0.6, 0.8} {
    \draw (\x,8) -- (\x,7.9) node[below, font=\small] {\x};
  }
  
  % Деления на оси Y (каждые 2 В)
  \foreach \y in {10, 15} {
    \draw (0,\y) -- (-0.01,\y) node[left, font=\small] {\y};
  }
  
  % Горизонтальные пунктирные линии
  \draw[red, dashed, thick] (0,15) -- (0.9,15);
  \node[red, right] at (0.75,15.5) {$u_C(\infty) = 15$ В};
  
  \draw[green!60!black, dashed, thick] (0,10) -- (0.9,10);
  \node[green!60!black, right] at (0.75,9.5) {$u_C(0_+) = 10$ В};
  
  % График функции u_C(t) = 15 - 5*exp(-6000*t)
  \draw[blue, thick, domain=0:0.9, samples=200, smooth] 
    plot (\x, {15 - 5*exp(-6000*\x/1000)});
  
  % Вертикальная пунктирная линия для τ
  \draw[orange, dashed] (0.167,8) -- (0.167,{15 - 5*exp(-1)});
  \node[orange, below] at (0.167,7.8) {$\tau$};
  
\end{tikzpicture}
\end{center}

График показывает переходный процесс заряда конденсатора от начального значения $10$ В до установившегося значения $15$ В с постоянной времени $\tau \approx 0{,}167$ мс.

\newpage

\section{Найти $i_2(t)$}

\begin{tabular}{l c}
  % Документация:
  % \begin{minipage}[t]{0.25\linewidth}
  % [t] — опция вертикального выравнивания: выравнивать верхнюю границу minipage по базовой линии окружения (top).
  % Альтернативы: [c] — по центру, [b] — по нижней границе.
  % {0.25\linewidth} — ширина minipage (25% от ширины строки).
  % Содержимое — строки в математическом режиме; '\\' делает перенос строки внутри minipage.
  % Символ '&' указывает на разделение колонок (minipage, вероятно, используется внутри табличной/выравнивающей среды).
  % \begin{tikzpicture}[scale=1,baseline=(current bounding box.center), european resistors]
  % — начало рисунка TikZ: scale задаёт масштаб, baseline выравнивает рисунок по указанной точке, 'european resistors' — стиль/библиотека для элементного оформления схем.
  \scalebox{1.2}{
    \begin{minipage}[c]{0.3\linewidth}
      \vspace{0.5cm}
      $E=10\,\text{В}$ \\ 
      $R_1=1\,\Omega$\\
      $R_2=1\,\Omega$ \\ 
      $R_3=1\,\Omega$ \\ 
      $R_4=1\,\Omega$ \\ 
      $L=1$\,\text{Гн} \\
    \end{minipage}
  } & \begin{tikzpicture}[scale=1,baseline=(current bounding box.center), european resistors, american inductors]
    \draw (0,-1) -- (0,0.5) to[L, l=$L$] (0,3) to[R, l=$R_1$] (0,5.5)
    to[short, -] (0,5.5) -- (2,5.5) to[short, -*] (2.5,5.5)
    -- (5.5,5.5) to[R, l=$R_4$] (5.5,-1) to[short, -*] (2.5,-1) -- (0,-1) -- (0,-1);
    \draw (2.5,5.5) to[short, i<=$i_2$] (2.5,4.5) to[R, l=$R_2$] (2.5,3) to[voltage source, l=$E$] (2.5,1) to[R, l=$R_3$] (2.5,-1);
    \draw (4,5.5) to[short, *-] (4,4.9);
    \draw (2.5,2.9) to[short, *-] (4,2.9) to[short, -*] (4,4.1) -- (4.1,4.9);
    \draw[fill=white] (4,4.9) circle (1.5pt);
    \draw[-{stealth}] ({4-0.3},{4.1+0.5}) arc (120:50:0.7);
  \end{tikzpicture} \\
\end{tabular}

\subsection{Решение}

\textbf{1. Начальное условие до коммутации (при $t < 0$, ключ замкнут)}

В начальный момент ключ замкнут, резистор $R_2$ закорочен, соотвественно ток на резисторе $R_2$ равен нулю $i_2 = 0$.
\[
i_2(0_-) = 0\,\text{А}
\]

I закон коммутации для данного тока нарушается.
\[
i_2(0_+) \neq i_2(0_-)
\]

\textbf{2. Установившийся режим после коммутации (при $t \to \infty$, ключ разомкнут)}

После размыкания ключа резистор $R_4$ включается в цепь. 
\[
i_{2 \text{ пр}_+} = \frac{E}{R_2 + R_3 + \frac{R_1R_4}{R_1 + R_4}} = \frac{10}{1 + 1 + \frac{1 \cdot 1}{1 + 1}} = \frac{10}{2 + \frac{1}{2}} = \frac{10}{\frac{5}{2}} = 4\,\text{А}
\]

\textbf{3. Постоянная времени $\tau$ и показатель $p$}

Для определения $\tau$ необходимо найти эквивалентное сопротивление относительно индуктивности при закороченном источнике ЭДС и разомкнутом ключе.

При закороченном источнике $E$ в ветви с $R_2$ и $R_3$:
\[
R_{\text{экв}} = R_1 + R_3 + \frac{R_2 \cdot R_4}{R_2 + R_4} = 1 + 1 + \frac{1 \cdot 1}{1 + 1} = 2 + \frac{1}{2} = \frac{5}{2} = 2{,}5\,\Omega
\]

Постоянная времени
\[
\tau = \frac{L}{R_{\text{экв}}} = \frac{1}{\frac{5}{2}} = \frac{2}{5} = 0.4\,\text{с}
\]

В свою очередь, показатель 
\[
p = -\frac{1}{\tau} = -\frac{5}{2} = -2.5\,\text{с}^{-1}
\]

\textbf{4. Свободная компонента}

Переходный процесс является апериодическим, так как в цепи присутствует только один энергонакопительный элемент. Чьи колебания описываются свободной компонентой тока.
\[
i_{2\text{ св}}(t) = A \cdot e^{pt}
\]

\textbf{5. Поиск постоянной интегрированния $A$}

Так как ток при разомкнутом ключе нам не известен, составим схему замещения для поиска $i_2(0_+)$.
Сначала определим ток через индуктивность до коммутации. При замкнутом ключе и в установившемся режиме $u_L = 0$, весь ток от источника $E$ идет через закороченный участок, следовательно
\[
i_L(0_-) = 0\,\text{А}
\]

Согласно I закону коммутации
\[
i_L(0_+) = i_L(0_-) = 0\,\text{А}
\]

Схема замещения в момент $t = 0_+$ (индуктивность заменена источником тока $i_L(0_+) = 0$ А):

\begin{center}
\begin{tikzpicture}[scale=1,baseline=(current bounding box.center), european resistors, american inductors]
    \draw (0,-1) -- (0,0.5) to[current source, l=$i_L(0_+)$] (0,3) to[R, l=$R_1$] (0,5.5)
    to[short, -] (0,5.5) -- (2,5.5) to[short, -*] (2.5,5.5)
    -- (5.5,5.5) to[R, l=$R_4$] (5.5,-1) to[short, -*] (2.5,-1) -- (0,-1) -- (0,-1);
    \draw (2.5,5.5) to[short, i<=$i_2$] (2.5,4.5) to[R, l=$R_2$] (2.5,3) to[voltage source, l=$E$] (2.5,1) to[R, l=$R_3$] (2.5,-1);
\end{tikzpicture}
\end{center}

Так как $i_L(0_+) = 0$ А, вся ветвь с $R_1$ фактически отключена (разрыв цепи). Ток $i_2(0_+)$ определяется только источником $E$ и резисторами $R_2$, $R_3$, $R_4$.

\[
i_2(0_+) = \frac{E}{R_2 + R_3 + R_4} = \frac{10}{1 + 1 + 1} = \frac{10}{3} \approx 3{,}33\,\text{А}
\]

Переходный процесс описывается уравнением
\[
\begin{aligned}
i_2(0_+) &= i_{2\text{ пр}}(0_+) + i_{2\text{ св}}(0_+) \\
\frac{10}{3} &= 4 + A \cdot e^{0} \\
A &= \frac{10}{3} - 4 = -\frac{2}{3} \approx -0.67\,\text{А}
\end{aligned}
\]

\textbf{6. Итоговое выражение для $i_2(t)$}
\[
i_2(t) = 4 - \frac{2}{3} \cdot e^{-2.5t}\,\text{А}
\]

\begin{center}
\begin{tikzpicture}[xscale=6, yscale=6]
  % Сетка
  \draw[gray!30, dashed, ystep=0.2, xstep=0.2] (0,3) grid (1.99,4.19);
  
  % Оси
  \draw[-{stealth}, thick] (0,3) -- (2,3) node[right] {$t$, с};
  \draw[-{stealth}, thick] (0,3) -- (0,4.2) node[above] {$i_2(t)$, А};
  
  % Деления на оси X (каждые 0.4 с)
  \foreach \x in {0.8, 1.2, 1.6} {
    \draw (\x,3) -- (\x,2.95) node[below, font=\small] {\x};
  }
  
  % Деления на оси Y
  \foreach \y in {3.33, 4} {
    \draw (0,\y) -- (-0.02,\y) node[left, font=\small] {\y};
  }
  
  % Горизонтальные пунктирные линии
  \draw[red, dashed, thick] (0,4) -- (1.8,4);
  \node[red, right] at (1.5,4.1) {$i_2(\infty) = 4$ А};
  
  \draw[green!60!black, dashed, thick] (0,3.33) -- (1.8,3.33);
  \node[green!60!black, right] at (1.5,3.23) {$i_2(0_+) = \frac{10}{3}$ А};
  
  % График функции i_2(t) = 4 - (2/3)*exp(-2.5*t)
  \draw[blue, thick, domain=0:1.8, samples=200, smooth] 
    plot (\x, {4 - (2/3)*exp(-2.5*\x)});
  
  % Вертикальная пунктирная линия для τ
  \draw[orange, dashed] (0.4,3) -- (0.4,{4 - (2/3)*exp(-1)});
  \node[orange, below] at (0.4,3) {$\tau$};
  
\end{tikzpicture}
\end{center}

График показывает переходный процесс тока $i_2$ от начального значения $\frac{10}{3}$ А до установившегося значения $4$ А с постоянной времени $\tau = 0{,}4$ с.

\section{Найти $i_2(t)$}

\begin{tabular}{l c}
  % Документация:
  % \begin{minipage}[t]{0.25\linewidth}
  % [t] — опция вертикального выравнивания: выравнивать верхнюю границу minipage по базовой линии окружения (top).
  % Альтернативы: [c] — по центру, [b] — по нижней границе.
  % {0.25\linewidth} — ширина minipage (25% от ширины строки).
  % Содержимое — строки в математическом режиме; '\\' делает перенос строки внутри minipage.
  % Символ '&' указывает на разделение колонок (minipage, вероятно, используется внутри табличной/выравнивающей среды).
  % \begin{tikzpicture}[scale=1,baseline=(current bounding box.center), european resistors]
  % — начало рисунка TikZ: scale задаёт масштаб, baseline выравнивает рисунок по указанной точке, 'european resistors' — стиль/библиотека для элементного оформления схем.
  \scalebox{1.2}{
    \begin{minipage}[c]{0.2\linewidth}
      \vspace{0.5cm}
      $E=1\,\text{В}$ \\
      $J=1\,\text{A}$ \\ 
      $R_2=1\,\Omega$ \\ 
      $R_3=1\,\Omega$ \\ 
      $C=1$\,\text{мкФ} \\
    \end{minipage}
  } & \begin{tikzpicture}[scale=1,baseline=(current bounding box.center), european resistors, american inductors]
    % Нижняя горизонтальная линия
    \draw (0,-1) -- (7.5,-1);
    
    % Ветвь Z_1: источник тока
    \draw (0,-1) to[current source, l=$J$] (0,3);
    
    % Ветвь Z_2: резистор R_2 с током i_2
    \draw (2.5,-1) to[short, i<=$i_2(t)$] (2.5,0.5) to[R, l=$R_2$] (2.5,2.5) to[short] (2.5,3);
    
    % Ветвь Z_3: источник напряжения (разрезан для ключа)
    \draw (5,-1) to[voltage source, l=$E$] (5,1) to[short] (5,2);
    
    % Ветвь Z_4: конденсатор
    \draw (7.5,-1) to[C, l=$C$] (7.5,3);
    
    % Верхняя горизонтальная линия
    \draw (0,3) -- (2.5,3);
    
    % Резистор между Z_3 и Z_4 сверху
    \draw (5,3) to[R, l=$R_3$] (7.5,3);
    
    % Переключающий ключ между Z_2 и Z_3
    % Общая точка на Z_2
    \draw (2.5,3) -- (4,3);
    
    % Верхняя точка контакта - прямое соединение к верху Z_3
    \draw[fill=black] (5,3) circle (1.5pt);
    
    % Нижняя точка контакта - через источник E
    \draw[fill=white] (5,2) circle (1.5pt);
    
    \draw (5,3) -- (4.9,2);
    \draw[fill=white] (4,3) circle (1.5pt);
    
    % Стрелка переключения
    \draw[-{stealth}] (5.2,2.4) arc (270:180:0.7);
  \end{tikzpicture} \\
\end{tabular}

\subsection{Решение}

\textbf{1. Начальное условие до коммутации (при $t < 0$, ключ замкнут)}

В начальный момент ключ замкнут, через резистор $R_2$ течёт ток от источника тока $J$. Следовательно,
\[
i_2(0_-) = J = 1\,\text{А}
\]

I закон коммутации для данного тока нарушается.
\[
i_2(0_+) \neq i_2(0_-)
\]

\textbf{2. Установившийся режим после коммутации (при $t \to \infty$, ключ разомкнут)}

После размыкания ключа резистор $R_3$ включается в цепь, но так как на конденсаторе $C$ разрыв цепи, ток через $R_3$ равен нулю, следовательно ток через $R_2$ определяется источником тока $J$.
\[
i_{2 \text{ пр}_+} = J = 1\,\text{А}
\]

\textbf{3. Постоянная времени $\tau$ и показатель $p$}

Для определения $\tau$ необходимо найти эквивалентное сопротивление относительно конденсатора при закороченном источнике ЭДС и разомкнутом ключе.
При закороченном источнике $E$ в ветви с $R_3$ сопротивление
\[
R_{\text{экв}} = R_2 + R_3 = 1 + 1 = 2\,\Omega
\]
Постоянная времени
\[
\tau = R_{\text{экв}} C = 2 \cdot 1\,\text{мкФ} = 2 \cdot 10^{-6}\,\text{с} = 2\,\text{мкс}
\]
В свою очередь, показатель
\[
p = -\frac{1}{\tau} = -\frac{1}{2 \cdot 10^{-6}} = -5 \cdot 10^{5}\,\text{с}^{-1}
\]
\textbf{4. Свободная компонента}

Переходный процесс является апериодическим, так как в цепи присутствует только один энергонакопительный элемент. Чьи колебания описываются свободной компонентой тока.
\[
i_{2\text{ св}}(t) = A \cdot e^{pt}
\]

\textbf{5. Поиск постоянной интегрированния $A$}

Так как ток при разомкнутом ключе нам не известен, составим схему замещения для поиска $i_2(0_+)$.
Сначала определим напряжение на конденсаторе до коммутации. При замкнутом ключе и в установившемся режиме $u_C = E$, следовательно,
\[
u_C(0_-) = E = 1\,\text{В}
\]

Согласно II закону коммутации
\[
u_C(0_+) = u_C(0_-) = 1\,\text{В}
\]

Схема замещения в момент $t = 0_+$ (конденсатор заменен источником напряжения $u_C(0_+) = 1$ В)
\begin{center}
\begin{tikzpicture}[scale=1,baseline=(current bounding box.center), european resistors, american inductors]
    % Нижняя горизонтальная линия
    \draw (0,-1) -- (7.5,-1);
    
    % Ветвь Z_1: источник тока
    \draw (0,-1) to[current source, l=$J$] (0,3);
    
    % Ветвь Z_2: резистор R_2 с током i_2
    \draw (2.5,-1) to[short, i<=$i_2(0_+)$] (2.5,0.5) to[R, l=$R_2$] (2.5,2.5) to[short] (2.5,3);
    
    % Ветвь Z_4: источник напряжения вместо конденсатора
    \draw (7.5,-1) to[voltage source, l=$u_C(0_+)$] (7.5,3);
    
    % Верхняя горизонтальная линия
    \draw (0,3) -- (2.5,3);
    
    % Резистор между Z_3 и Z_4 сверху
    \draw (2.5,3) to[R, l=$R_3$] (7.5,3);
\end{tikzpicture}
\end{center}

Так как конденсатор заменен источником напряжения $u_C(0_+) = 1$ В, то его теперь также необходимо учитывать при расчёте тока $i_2(0_+)$.
\[
i_2(0_+) = J + \frac{u_C(0_+)}{R_2 + R_3} = 1 + \frac{1}{1 + 1} = 1 + \frac{1}{2} = 1.5\,\text{А}
\]

Переходный процесс описывается уравнением
\[
\begin{aligned}
i_2(0_+) &= i_{2\text{ пр}}(0_+) + i_{2\text{ св}}(0_+) \\
1.5 &= 1 + A \cdot e^{0} \\
A &= 1.5 - 1 = 0.5\,\text{А}
\end{aligned}
\]

\textbf{6. Итоговое выражение для $i_2(t)$}

\[
i_2(t) = 1 + 0.5 \cdot e^{-5 \cdot 10^{5}t}
\]

\begin{center}
\begin{tikzpicture}[xscale=1.5, yscale=4]
  % Сетка (t от 0 до 10 мкс, i от 0 до 1.6 А в координатах рисунка)
  \draw[gray!30, dashed, ystep=0.1, xstep=0.5] (0,0) grid (10.49,1.69);
  
  % Оси
  \draw[-{stealth}, thick] (0,0) -- (10.5,0) node[right] {$t$, мкс};
  \draw[-{stealth}, thick] (0,0) -- (0,1.7) node[above] {$i_2(t)$, А};
  
  % Деления на оси X (каждые 2 мкс)
  \foreach \x in {2,4,6,8,10} {
    \draw (\x,0) -- (\x,-0.02) node[below, font=\small] {\x};
  }
  
  % Деления на оси Y (0, 1, 1.5 А)
  \foreach \y/\lbl in {0/0,1/1,1.5/1.5} {
    \draw (0,\y) -- (-0.02,\y) node[left, font=\small] {\lbl};
  }
  
  % Горизонтальные пунктирные линии
  \draw[red, dashed, thick] (0,1) -- (10,1);
  \node[red, right] at (8.5,1.1) {$i_2(\infty) = 1$ А};
  
  \draw[green!60!black, dashed, thick] (0,1.5) -- (10,1.5);
  \node[green!60!black, right] at (8.5,1.6) {$i_2(0_+) = 1.5$ А};
  
  % График функции i_2(t) для t в мкс: exp(-5e5 * t[с]) = exp(-5e5 * t[мкс]*1e-6) = exp(-0.5 * t[мкс])
  \draw[blue, thick, domain=0:10, samples=600, smooth] 
    plot (\x, {1 + 0.5*exp(-0.5*\x)});
  
  % Вертикальная пунктирная линия для τ (τ = 2 мкс)
  \draw[orange, dashed] (2,0) -- (2,{1 + 0.5*exp(-1)});
  \node[orange, below] at (2,0) {$\tau$};
  
\end{tikzpicture}
\end{center}

График показывает переходный процесс тока $i_2$ от начального значения $1.5$ А до установившегося значения $1$ А с постоянной времени $\tau = 2$ мкс.


\end{document}